% Generated draft: requires mathematical and source review.
\documentclass[11pt]{article}
\usepackage[margin=25mm]{geometry}
\usepackage{fontspec}
\setmainfont{Latin Modern Roman}
\usepackage{amsmath,amssymb,mathtools,unicode-math}
\setmathfont{Latin Modern Math}
\usepackage{enumitem,needspace,xurl,hyperref}
\setlength{\parindent}{0pt}
\setlength{\parskip}{5pt}
\setlength{\emergencystretch}{3em}
\newcommand{\N}{\mathbb{N}}
\newcommand{\Z}{\mathbb{Z}}
\newcommand{\Q}{\mathbb{Q}}
\newcommand{\R}{\mathbb{R}}
\newcommand{\C}{\mathbb{C}}
\newcommand{\F}{\mathbb{F}}
\newcommand{\A}{\mathbb{A}}
\newcommand{\PP}{\mathbb{P}}
\newcommand{\E}{\mathbb{E}}
\newcommand{\eps}{\varepsilon}
\newcommand{\dd}{\,\mathrm{d}}
\begin{document}
\section*{UnsolvedMath ID: 700768}
\textbf{Source problem code:} catalogue entry
\subsection{Definitions and mathematical statement}
A finite semifield, in the usual coordinatizing convention, is a finite set \(S\) with operations \(+\) and \(\cdot\) such that \((S,+)\) is an abelian group with identity \(0\), multiplication has an identity \(1\ne0\), both distributive laws hold, and, for every \(a\ne0\) and every \(b\in S\), the equations \(a\cdot x=b\) and \(y\cdot a=b\) each have a unique solution. Multiplication need not be associative. The source does not explicitly specify its semifield convention; this definition states the conventional interpretation used in this draft.

A projective plane is non-Desarguesian if Desargues' theorem does not hold in it. A collineation of a projective plane \(\Pi\) consists of bijections on its points and on its lines that preserve incidence. Let \(\operatorname{Col}(\Pi)\) denote the group of all such collineations. A group \(G\) is soluble if its derived series, defined by \(G^{(0)}=G\) and \(G^{(i+1)}=[G^{(i)},G^{(i)}]\), satisfies \(G^{(n)}=\{1\}\) for some integer \(n\ge0\).

\textbf{Question.} For every finite non-Desarguesian projective plane \(\Pi\) coordinatizable over a semifield, is \(\operatorname{Col}(\Pi)\) soluble? Equivalently, does every such \(\Pi\) satisfy
\[
\exists n\in\Z_{\ge0}\qquad \operatorname{Col}(\Pi)^{(n)}=\{1\}?
\]
The source qualifies the hypothesis on \(\Pi\) by saying that the corresponding regular set is closed under addition, and refers to Archive problem 10.48 for that terminology. That referenced definition is not included in the supplied record, so its precise conventions and correspondence with semifield coordinatization require checking against the cited source. The full collineation group is used here because the supplied statement specifies no subgroup or stabilizer.
\subsection{Short English statement}
Must the full collineation group of every finite non-Desarguesian projective plane defined over a semifield be soluble? The source describes the plane hypothesis through an associated regular set closed under addition; its reference to Archive problem 10.48 requires source review.
\subsection{Sources}
\begin{itemize}
\item Local UnsolvedMath registry, supplied complete record with permanent ID \(700768\), exact source problem code \texttt{catalogue entry}, title \emph{Kourovka Notebook Problem 11.76}; upstream revision \texttt{d371ee917ac9647fbb06cd6fc529deec6f3dda82}. The supplied registry status is \texttt{open}; this label is not an independently verified mathematical or literature assessment.
\item \emph{The Kourovka Notebook (Unsolved Problems in Group Theory)}, 11th issue (1990), problem 11.76; also identified in the supplied record with the 21st issue edition (2026), \texttt{arXiv:1401.0300v48}. The supplied collection description names E. I. Khukhro and V. D. Mazurov as editors. These bibliographic details are supplied attribution and have not been independently checked.
\item Proposed by N. D. Podufalov, 1990. The source labels this a well-known problem.
\item The source explicitly refers to Archive problem 10.48 for the corresponding regular set and states that the plane hypothesis means this set is closed under addition. The referenced archive entry is not supplied.
\end{itemize}
\end{document}
